% =============================================================================
%  chart_mpi.tex  —  Diagramme Speedup / Efficacité MPI
% =============================================================================
% Auteur     : Merbouche Abdelatif
% Module     : Calculs Parallèles — M1 HPC, USTHB
% Encadrante : Mme Baba Ali
% Année      : 2025/2026
%
% Description :
%   Reproduction vectorielle du diagramme de performance généré par
%   l'application web du projet MPI. Ce graphique présente l'évolution
%   du speedup et de l'efficacité de l'algorithme k-NN sur un dataset
%   de classification, selon le nombre de processus MPI.
%
% Compilation :
%   pdflatex chart_mpi.tex
%   (ou lualatex chart_mpi.tex)
%
% Résultat : chart_mpi.pdf (graphique vectoriel haute qualité)
% =============================================================================

\documentclass[tikz,border=5pt]{standalone}
\usepackage[utf8]{inputenc}
\usepackage[T1]{fontenc}
\usepackage[french]{babel}
\usepackage{pgfplots}
\pgfplotsset{compat=1.18}
\usepackage{tikz}
\usetikzlibrary{positioning}

% --- Palette de couleurs du projet ---
\definecolor{mpiBlue}{HTML}{2563EB}
\definecolor{mpiGreen}{HTML}{059669}
\definecolor{mpiGray}{HTML}{94A3B8}
\definecolor{mpiDark}{HTML}{1E3A5F}

\begin{document}
\begin{tikzpicture}

% ============================================================================
%  DONNÉES EXPÉRIMENTALES
%  k-NN — Dataset classification (similaire Iris)
%  Ts = 0.025 s (temps séquentiel de référence)
% ============================================================================
\pgfplotstableread{
  p    speedup    efficiency
  1    1.00       100.0
  2    1.89        94.7
  4    3.47        86.8
  6    4.90        81.7
  8    6.25        78.1
  10   7.58        75.8
  12   8.62        71.8
  16  10.87        67.9
  20  12.50        62.5
}\datatable

% ============================================================================
%  AXE Y GAUCHE — SPEEDUP
% ============================================================================
\begin{axis}[
    width=14cm,
    height=8cm,
    axis y line*=left,
    xlabel={\textbf{Nombre de processus MPI ($p$)}},
    ylabel={\textbf{Speedup $S = T_s / T_p$}},
    xmin=0, xmax=21,
    ymin=0, ymax=21,
    xtick={1,2,4,6,8,10,12,16,20},
    ytick={0,2,4,6,8,10,12,14,16,18,20},
    grid=major,
    grid style={dashed, mpiGray!40},
    minor y tick num=1,
    legend style={
      at={(0.02,0.98)},
      anchor=north west,
      draw=mpiGray,
      fill=white,
      font=\small
    },
    title={\Large\textbf{Speedup et efficacite --- k-NN (dataset classification)}},
    title style={yshift=0.3cm, text=mpiDark},
    label style={font=\small},
    tick label style={font=\small},
    axis line style={mpiDark},
]

% Speedup observé (courbe + marqueurs)
\addplot[
  color=mpiBlue,
  mark=*,
  mark size=2.5pt,
  mark options={fill=white, draw=mpiBlue, line width=1.2pt},
  line width=2pt,
  smooth,
  tension=0.2
]
  table [x=p, y=speedup] {\datatable};
\addlegendentry{Speedup observé}

% Speedup idéal (linéaire S = p)
\addplot[
  color=mpiGray,
  dashed,
  line width=1.2pt,
  domain=0:21,
  samples=50
]
  {x};
\addlegendentry{Speedup idéal ($S = p$)}

% Annotation du dernier point
\node[anchor=south west, font=\footnotesize\bfseries, color=mpiBlue]
  at (axis cs:18,12.5) {$12{,}50\times$};
\draw[mpiGray, ->, thick]
  (axis cs:17,11.8) -- (axis cs:19.5,12.5);

\end{axis}

% ============================================================================
%  AXE Y DROITE — EFFICACITÉ (%)
% ============================================================================
\begin{axis}[
    width=14cm,
    height=8cm,
    axis y line*=right,
    axis x line=none,
    ylabel={\textbf{Efficacité $E = S/p$ (\%)}},
    ymin=0, ymax=110,
    ytick={0,20,40,60,80,100},
    legend style={
      at={(0.02,0.72)},
      anchor=north west,
      draw=mpiGray,
      fill=white,
      font=\small
    },
    label style={font=\small},
    tick label style={font=\small},
    axis line style={mpiDark},
]

% Efficacité observée (courbe + marqueurs carrés)
\addplot[
  color=mpiGreen,
  mark=square*,
  mark size=2.5pt,
  mark options={fill=white, draw=mpiGreen, line width=1.2pt},
  line width=2pt,
  smooth,
  tension=0.2
]
  table [x=p, y=efficiency] {\datatable};
\addlegendentry{Efficacité (\%)}

% Efficacité idéale (100 %)
\addplot[
  color=mpiGray,
  dashed,
  line width=1.2pt,
  domain=0:21,
  samples=50
]
  {100};
\addlegendentry{Efficacité idéale (100 \%)}

% Annotation du dernier point
\node[anchor=north east, font=\footnotesize\bfseries, color=mpiGreen]
  at (axis cs:18.5,62.5) {$62{,}5\,\%$};
\draw[mpiGray, ->, thick]
  (axis cs:17,55) -- (axis cs:19.5,62);

\end{axis}

\end{tikzpicture}
\end{document}
